\documentclass[10pt,a4paper]{article}

\usepackage[margin=1.5cm,top=1.3cm,bottom=1.3cm]{geometry}
\usepackage[T1]{fontenc}
\usepackage[utf8]{inputenc}
\usepackage{lmodern}
\usepackage{microtype}
\usepackage{booktabs}
\usepackage{tabularx}
\usepackage{enumitem}
\usepackage{xcolor}
\usepackage[hidelinks]{hyperref}

\setlist{nosep,leftmargin=1.2em}
\setlength{\parindent}{0pt}
\setlength{\parskip}{3pt}
\pagestyle{empty}
\newcommand{\src}[1]{{\color{black!55}\small[#1]}}
\newcommand{\sect}[1]{\vspace{5pt}{\large\bfseries #1}\par\vspace{1pt}}

\begin{document}

{\LARGE\bfseries Week 1 summary: sources and salient points}\par
\vspace{2pt}
\textbf{Satchel}, a scaled-down clone of Dropbox Magic Pocket.
Team: Parshv Joshi (@frectrium), @VyasGuru, @JasMehta08.
Source IDs refer to the Source Log (\texttt{docs/sources/README.md}).

\sect{Sources found}
\begin{tabularx}{\linewidth}{@{}l r X@{}}
\toprule
Type & Count & Most important \\
\midrule
Dropbox Tech Blog (B, primary) & 15 & \emph{Inside the Magic Pocket} [B2], \emph{Scaling to exabytes and beyond} [B1], \emph{Pocket watch} [B3] \\
Conference talks (T, primary) & 10 & QCon SF 2022 [T8], SNIA SDC 2024 [T9], Storage Field Day 2018 [T5] \\
SEC filings (F, primary) & 2 & 2018 S-1 [F1], 2018 10-K [F2] \\
Media (M, context only) & 8 & GeekWire [M5], Data Center Knowledge [M6] \\
Cloud pricing (P) & 9 & S3 pricing, ap-south-1 [P1]; Azure Retail Prices API, Central India [P5] \\
On-prem prices (H) & 13 & HDD \$/TB [H2], PUE [H6], Indian colocation [H8], IP transit [H9] \\
\bottomrule
\end{tabularx}

\sect{Problem: why Dropbox left S3}
By 2016 Dropbox held 500\,PB of user data, up more than 12x from about 40\,PB in 2012. File content lived on S3, while metadata ran in Dropbox's own data centers \src{B1}. Dropbox gave two reasons for building its own storage: \textbf{performance}, from controlling the stack end to end, and \textbf{better unit economics} at its scale \src{B1}. A later talk called cost the leading reason \src{T6}. An expiring AWS storage commitment set the migration deadline \src{T4}.

\sect{Solution: Magic Pocket}
\begin{itemize}
  \item \textbf{Data model:} an immutable block store. Blocks are up to 4\,MB and keyed by SHA-256, each stored in at least two US zones and replicated within each zone \src{B2}.
  \item \textbf{Indexes and cells:} a sharded MySQL Block Index maps each block to a cell and bucket. Each cell has storage nodes (OSDs), plus a Master and Volume Managers that run repair and erasure coding outside the request path \src{B2}.
  \item \textbf{Encoding:} new data is replicated (4 copies per zone), then erasure coded. The 2024 production schemes are RS(6,3) and LRC(12,2,2) \src{T8, T9}.
  \item \textbf{Verification and hardware:} six continuous verifiers check durability and use over half the disk and database load \src{B3}. Storage runs on dense SMR servers, and the storage nodes moved from Go to Rust for memory control \src{B4}.
\end{itemize}

\sect{Outcome}
\begin{itemize}
  \item \textbf{Migration:} over 90\% of user data was served in-house by 2015-10-07, ``without any major service disruptions or any loss of data'' \src{B1}. The duplicated copies were retired in Q4 2016 \src{F1}.
  \item \textbf{Cost:} infrastructure costs fell \textbf{\$39.5M} in 2016 and \textbf{\$35.1M} in 2017, and gross margin rose from 33\% (2015) to 67\% (2017) \src{F1}. The S-1 groups the migration with account cleanup under ``Infrastructure Optimization'', so these figures are not Magic Pocket savings alone.
  \item \textbf{Hybrid result:} AWS still serves the remaining data, delivery and European data residency \src{F1, T8}.
  \item \textbf{Later savings:} SMR drives saved more than 20\% over the previous storage design \src{B7}, and the cold tier cut disk use by 25\% \src{B5}.
\end{itemize}

\sect{Inputs for the cost exercise so far}
\begin{tabularx}{\linewidth}{@{}X X X X@{}}
\toprule
List prices & Storage per GB-month & Writes per 1,000 & Internet egress per GB \\
\midrule
S3 Standard, ap-south-1 \src{P1} & \$0.025 down to \$0.023 & \$0.005 & \$0.1093 down to \$0.08 \\
Azure Hot ZRS, Central India \src{P5, P6} & \$0.025 down to \$0.023 & \$0.006875 & \$0.12 down to \$0.08 \\
\bottomrule
\end{tabularx}
\par\vspace{3pt}
On-prem references: enterprise HDDs at \$27.92--44.17 per TB in August 2026, up 15--48\% since March \src{H2}; PUE 1.52 in Asia-Pacific \src{H6}; Indian colocation at Rs 25,000--60,000 per rack per month plus power \src{H8}; IP transit from \$0.03 per Mbps per month \src{H9}; servers depreciated over 5--6 years \src{H13}.

\sect{Open questions for Week 2}
\begin{enumerate}
  \item How long does Dropbox keep data replicated before erasure coding it? No source says, and it decides Satchel's parameters (\#14).
  \item Which workload size and year do we model: 500\,PB \src{B1} or about 600\,PB \src{B14} (\#15)?
  \item How do we update the 2019 full-server cost \src{H3} to 2026, given volatile drive prices \src{H2}?
  \item Which S3 and Azure tiers match Dropbox's mostly cold, instant-read workload (\#15)?
\end{enumerate}

\end{document}
